%! TeX root = ../main.tex

\addchap*{Abstract}

In this thesis we present an impurity solver for dynamical mean-field theory.
The use of a restricted active basis set allows us to couple the impurity to hundreds of bath sites
at polynomial cost.
Dynamical quantities are computed directly on the real axis
without the introduction of any broadening.
In contrast to most previous works the self-energy is calculated from four correlators
and its imaginary part is negative by construction.
We study the Mott metal-to-insulator transition on the Bethe lattice at zero temperature
and observe sharp features at the inner edge of the Hubbard bands.
The quasiparticle weight and impurity double occupation
are compared against results of numerical renormalization group calculations.

\selectlanguage{ngerman}

\addchap*{Kurzfassung}

In dieser Arbeit präsentieren wir einen Algorithmus zur Lösung des \texten{Anderson}'schen
Störstellenprolems für dynamische Molekularfeldtheorie.
Durch den Gebrauch eines eingeschränkten Hilbertraums kann diese Störstelle in polynomialer Zeit
an hunderte Badgitterplätze gekoppelt werden.
Dynamische Größen werden direkt an der reellen Achse ohne Verbreiterung berechnet.
Im Gegensatz zu vielen vorigen Arbeiten
wird die Selbstenergie mittels vier Korrelatoren berechnet
und ihr Imaginärteil ist per Konstruktion negative.
Wir analysieren den \texten{Mott} Übergang von einem Metall zu Isolator
auf dem \texten{Bethe} Gitter am absoluten Temperaturnullpunkt
und beobachten schmale Strukturen auf den inneren Seiten der Hubbardbänder.
Das Gewicht des Quasiteilchen und die Doppelbesetzung der Störstelle werden
gegen Resultate der numerische Renormierungsgruppe verglichen.

\selectlanguage{english}
